\documentclass[11pt,twoside]{article}
\newcommand{\thang}[1]{{\color{blue}\small thang:#1}}

\usepackage{amssymb}
\usepackage{amsmath}
\usepackage{latexsym}
\usepackage{color}
\usepackage{graphics}
\usepackage{xspace}

% Commands for special characters
\newcommand{\mybackslash}{\char'134}
\newcommand{\myleftbracket}{\char'173}
\newcommand{\myrightbracket}{\char'175}
\newcommand{\mypercent}{\char'045}
\newcommand{\myunderscore}{\char'137}

% The 'ifthen' package supports Boolean flags
\usepackage{ifthen}
% The 'solutions' flag determines whether this is the original handout
%    or a solution
\newboolean{solutions}
\setboolean{solutions}{False}  % Default is original handout

% Uncomment the next line before starting on the solutions
% \setboolean{solutions}{True}

\newcommand{\coursenumber}{CS 4104}
\newcommand{\coursetitle}{Data and Algorithm Analysis}
\newcommand{\docdate}{October 02, 2023}
\newcommand{\duedate}{October 13, 2023}
\newcommand{\homeworknumber}{3}

% The document title depends on whether these are solutions
\ifthenelse{\boolean{solutions}}{% Then
\newcommand{\doctitle}{Solutions to Homework Assignment \homeworknumber}
}{% Else
\newcommand{\doctitle}{Homework Assignment \homeworknumber}
}

% The document date depends on whether these are solutions
\ifthenelse{\boolean{solutions}}{% Then
\renewcommand{\docdate}{\duedate}
}{% Else
}

% If you are the author, so put your name here
\renewcommand{\author}{Thang Hoang}

\pagestyle{myheadings}
\markboth{\hfill\doctitle\hfill\docdate}{\docdate\hfill\doctitle\hfill}

\addtolength{\textwidth}{1.00in}
\addtolength{\textheight}{1.00in}
\addtolength{\evensidemargin}{-1.00in}
\addtolength{\oddsidemargin}{-0.00in}
\addtolength{\topmargin}{-.50in}
\setlength{\footskip}{0pt}

\newcommand{\polyreduce}{\leq_{\mathrm{P}}}

\newcounter{problem}
\setcounter{problem}{0}
\newcommand{\problem}[1]{%
\refstepcounter{problem}\noindent\textbf{[#1] \arabic{problem}.}}

\newcommand{\solution}{\bigskip\hrule\bigskip}
\newcommand{\problembreak}{\bigskip\hrule\bigskip}

\renewcommand{\theenumi}{\textbf{\Alph{enumi}}}
\renewcommand{\theenumii}{\textbf{\roman{enumii}}}

\newcommand{\emptysequence}{\Lambda}
\newcommand{\emptystring}{\lambda}

% Pseudocode comment symbol
\newcommand{\comment}{\textbf{//}\ \ }

% Pseudocode line numbering
\newcounter{pseudocode}
\newcommand{\firstline}{\setcounter{pseudocode}{0}\linenumber}
\newcommand{\linenumber}{\refstepcounter{pseudocode}\thepseudocode}
\newcommand{\pseudoindent}{\hspace*{26pt}}

\newcommand{\nil}{\mbox{\textsc{nil}}}

% Mathematical symbols
\newcommand{\grid}{\mathcal{G}}  % Grid graph
\newcommand{\integers}{\mathbb{Z}}  % Integers
\newcommand{\naturals}{\mathbb{N}}  % Natural numbers
\newcommand{\reals}{\mathbb{R}}  % Real numbers

% Probability
\newcommand{\expect}[1]{\mathbf{E}\left[#1\right]}
\newcommand{\prob}[1]{\mathrm{Pr}\left[#1\right]}
\newcommand{\var}[1]{\mathrm{Var}\left[#1\right]}

% Logic
\newcommand{\NOT}[1]{\neg #1}
\newcommand{\AND}{\wedge}
\newcommand{\OR}{\vee}
\newcommand{\clause}[1]{\left(#1\right)}

\newlength{\problemoffset}
\setlength{\problemoffset}{0.5in}

% Decision problem macro
% A command for formatting decision problems a la Garey and Johnson
\newcommand{\decision}[3]{%     \decision{NAME}{INSTANCE}{QUESTION}
\begin{list}{}{
\setlength{\leftmargin}{\problemoffset}
\setlength{\rightmargin}{\problemoffset}
\setlength{\parsep}{0pt}
\setlength{\itemsep}{2pt}
\setlength{\topsep}{\itemsep}
\setlength{\partopsep}{\itemsep}
}
\item
{\textsc{#1}}
\item
{INSTANCE: #2}
\item
{QUESTION: #3}
\end{list}
}

% Optimization problem macro
\newcommand{\optimization}[3]{%  \optimization{NAME}{INSTANCE}{SOLUTION}
\begin{list}{}{
\setlength{\leftmargin}{\problemoffset}
\setlength{\rightmargin}{\problemoffset}
\setlength{\parsep}{0pt}
\setlength{\itemsep}{2pt}
\setlength{\topsep}{\itemsep}
\setlength{\partopsep}{\itemsep}
}
\item
{\rule{0pt}{14pt}\textsc{#1}}
\item
{INSTANCE: #2}
\item
{SOLUTION: #3}
\end{list}
}

\newcommand{\reaches}{\leadsto}

% Table layout
\newcommand{\toprule}{\rule[11pt]{0pt}{2pt}}
\newcommand{\bottomrule}{\rule[-6pt]{0pt}{0pt}}
\newenvironment{raggedpars}[1][2.0in]{%
\begin{minipage}[t]{#1}\raggedright\toprule}%
{\bottomrule\end{minipage}}

% Dynamic programming macros
\newlength{\arrowwidth}
\setbox3=\hbox{$\nwarrow$}
\setlength{\arrowwidth}{\wd3}
\newcommand{\optnwarrow}[1]{\if1#1\nwarrow%
\else\rule{\arrowwidth}{0pt}\fi}
\newcommand{\optuparrow}[1]{\if1#1\uparrow%
\else\rule{\arrowwidth}{0pt}\fi}
\newcommand{\optleftarrow}[1]{\if1#1\leftarrow%
\else\rule{\arrowwidth}{0pt}\fi}
% Use \tablebox to specify any combination of arrows, plus the value
\newcommand{\tablebox}[4]{%
\setlength{\arraycolsep}{0.0pt}%
\begin{array}{cc}%
\optnwarrow{#1} & \optuparrow{#2} \\%
\optleftarrow{#3} & #4%
\end{array}%
}
% Use \tableboxred to specify any combination of arrows, plus the value
% The value will be red; arrows are made red if 2 used instead of 1
\newcommand{\optnwarrowred}[1]{\if1#1\nwarrow%
\else\if2#1\textcolor{red}{\nwarrow}\else\rule{\arrowwidth}{0pt}\fi\fi}
\newcommand{\optuparrowred}[1]{\if1#1\uparrow%
\else\if2#1\textcolor{red}{\uparrow}\else\rule{\arrowwidth}{0pt}\fi\fi}
\newcommand{\optleftarrowred}[1]{\if1#1\leftarrow%
\else\if2#1\textcolor{red}{\leftarrow}\else\rule{\arrowwidth}{0pt}\fi\fi}
\newcommand{\tableboxred}[4]{%
\setlength{\arraycolsep}{0.0pt}%
\begin{array}{cc}%
\optnwarrowred{#1} &%
\optuparrowred{#2} \\%
\optleftarrowred{#3} &%
\textcolor{red}{#4}%
\end{array}%
}

% Allow really sloppy paragraphs that do not generate overfull or
%    underfull hbox's
\newenvironment{SLOPPY}{\begin{sloppypar}\hbadness=10000}{\end{sloppypar}}

% Definitions for this homework
\newcommand{\extent}[1]{\mathrm{extent}(#1)}
\newcommand{\opt}[2]{\mbox{\textsc{opt}}(#1,#2)}
\newcommand{\gap}{\mbox{\texttt{-}}}
\newcommand{\rewriterule}[2]{#1\to #2}
\newcommand{\rewrites}[2][]{\mathop{\Longrightarrow}\limits_{#1}^{#2}}
\newcommand{\reduces}{\leq}
\newcommand{\polyreduces}{\leq_P}
\newcommand{\mathsc}[1]{\mbox{\normalfont\textsc{#1}}}
\newcommand{\NP}{\mathcal{NP}}
\renewcommand{\P}{\mathcal{P}}
\newcommand{\describes}{\vdash}

\begin{document}

\thispagestyle{empty}

\begin{center}
\begin{tabular}{lcr}
\multicolumn{3}{c}{\Large\textbf{\coursenumber}}
\\[0.04in]
\multicolumn{3}{c}{\Large\textbf{\doctitle}}
% If these are solutions, then include the author's (student's) name
\ifthenelse{\boolean{solutions}}{% Then
\\[0.04in]\multicolumn{3}{c}{\large\textbf{\author}}
}{} % Else omit
% If these are solutions, then the date is the due date
\ifthenelse{\boolean{solutions}}{% Then
\\[0.10in]\multicolumn{3}{c}{\duedate}
}{% Else, put given and due dates
\\[0.10in]
\textbf{Given:} \docdate
& \hspace*{1.0in} &
\textbf{Due:} \duedate
}
\end{tabular}
\end{center}

% If these are solutions, then we do not include the directions
\ifthenelse{\boolean{solutions}}{} % No directions
{
% Original document includes directions
\begingroup % This allows an argument that contains multiple paragraphs
\paragraph{General directions.}
The point value of each problem is shown in [ ].
Each solution must include all details and
an explanation of why the given solution is correct.
\textbf{\textcolor{red}{In particular,
write complete sentences.
A correct answer without an explanation is worth no credit.}}
The completed assignment must be submitted on Canvas as a PDF by 5:00 PM
on \duedate.
\textbf{No late homework will be accepted.}

\paragraph{Digital preparation of your solutions is mandatory.
This includes digital preparation of any drawings.}
Use of \LaTeX\ is required.
Also,
\textbf{please include your name.}

\paragraph{Use of \LaTeX\ (required).}
\begin{itemize}
\item Retrieve this \LaTeX\ source file,
named
\texttt{homework\homeworknumber.tex},
from the course web site.
\item Rename the file
\textit{$<$Your VT PID$>$}\texttt{{\myunderscore}solvehw\homeworknumber.tex},
For example,
for the instructor,
the file name would be
\texttt{hoang{\myunderscore}solvehw\homeworknumber.tex}.

\item
Use a \textbf{text editor}
(such as \texttt{vi}, \texttt{emacs}, or \texttt{pico})
to accomplish the next three steps.
Alternately,
use Overleaf as your \LaTeX\ platform.

\item
Uncomment the line

\texttt{{\mypercent} 
{\mybackslash}setboolean{\myleftbracket}solutions{\myrightbracket}%
{\myleftbracket}True{\myrightbracket}}

in the document preamble by deleting the \texttt{\mypercent}.

\item
Find the line

\texttt{{\mybackslash}renewcommand%
{\myleftbracket}{\mybackslash}author{\myrightbracket}%
{\myleftbracket}Thang Hoang{\myrightbracket}}

and replace the instructor's name with your name.

\item
Enter your solutions where you find
the \LaTeX\ comments

\texttt{{\mypercent} PUT YOUR SOLUTION HERE}

\item
Generate a PDF and turn it in on Canvas by 5:00 PM on \duedate.
\end{itemize}
\endgroup

\problembreak

\newpage

}

\problem{30}
Here, we generalize the Longest Common Subsequence (LCS) problem
discussed in class
to a new problem that can also be solved by dynamic programming.

Let $\Sigma$ be any finite alphabet containing at least two symbols.
A \textit{balloon} in $\Sigma$ is a subset of $\Sigma$
containing exactly two symbols.
For example,
if $\Sigma=\{a,c,g,t\}$,
then $\{c,t\}$ and $\{a,t\}$ are balloons in $\Sigma$,
while $\emptyset$, $\{a,c,g\}$, and $\{t\}$ are not.
Note that there are $\binom{|\Sigma|}{2}$ balloons in any $\Sigma$.

A \textit{party} over $\Sigma$ is a sequence of balloons in $\Sigma$;
the empty party is denoted $\emptysequence$.
For example,
if $\Sigma=\{a,c,g,t\}$,
\begin{displaymath}
\{c,t\},\{a,t\},\{a,g\}
\end{displaymath}
is a party over $\Sigma$ of length 3.
A party $B=B_1,B_2,\ldots,B_n$ over $\Sigma$ of length $n$
\textit{describes} a string $A=\alpha_1,\alpha_2,\ldots,\alpha_n$
over $\Sigma$ of length $n$
if $\alpha_i\in B_i$, for $1\leq i\leq n$.
This is written $B\describes A$.
For example,
\begin{eqnarray*}
\{c,t\},\{a,t\},\{a,g\} & \describes & tta
\end{eqnarray*}
and
\begin{eqnarray*}
\{c,t\},\{a,t\},\{a,g\} & \describes & ctg.
\end{eqnarray*}

The optimization problem to be solved is:
\optimization{Balloon Longest Common Subsequence (BLCS)}%
{Parties $X=X_1,X_2,\cdots,X_m$
and $Y=Y_1,Y_2,\cdots,Y_n$
over an alphabet $\Sigma$.}%
{Strings $A=\alpha_1\alpha_2\cdots\alpha_m$ and $B=\beta_1\beta_2\cdots\beta_n$
such that $X\describes A$, $Y\describes B$,
and the length of an LCS of $A$ and $B$ is maximum.}

In the subsequent steps,
you are to develop a dynamic programming solution to the problem
and illustrate your solution with a particular example.

{\bfseries
\begin{enumerate}

\item
For an instance consisting of two parties $X$ and $Y$,
explain carefully what subinstances
you identify for solution
and argue that the problem possesses optimal substructure.

\item
Develop a recurrence for the optimal values $c(i,j)$
for all the subinstances identified.
Explain your recurrence.

\item
Let $\Sigma=\{a,c,g,t\}$.
Write a small program in your favorite programming language
to generate a random party over $\Sigma$ of length $n$.
Include the code for your program in your solution.
Generate an example random party $X=X_1,X_2,X_3,X_4,X_5,X_6,X_7$
of length 7
and an example random party $Y=Y_1,Y_2,Y_3,Y_4,Y_5$
of length 5.

\item
Figure~\ref{figure:emptyBLCS}
gives the template for the table
to contain the results of dynamic programming
on the BLCS instance $(X,Y)$ from Item C.
Fill in the table with the $c(i,j)$ values
obtained from your recurrence.
Include the arrows\footnote{%
If you use \LaTeX,
then the \texttt{{\char'134}tablebox} macro in this document
is capable of inserting the arrows.
},
as was done in the example of LCS in class.
What is the optimal objective value for the instance $(X,Y)$?

\item
Use backtracing to obtain $A$ and $B$
such that $X\describes A$, $Y\describes B$,
and $A$ and $B$ have a longest possible LCS.
Especially show the arrows\footnote{%
If you use \LaTeX,
then the \texttt{{\char'134}tableboxred} macro in this document is
capable of typesetting the relevant values and arrows
in \textcolor{red}{red}.
}
in the table that are
part of finding $A$ and $B$.
What are the $A$ and $B$ you obtain?

\end{enumerate}
}

\solution

% PUT YOUR SOLUTION HERE

\problembreak
\begin{figure}[tp]
\begin{displaymath}
\begin{array}{rc|c|c|c|c|c|c|c|}
\multicolumn{2}{c}{} &
\multicolumn{6}{c}{j}
\\[6pt]
& c(i,j) &
\begin{array}{c}\\ 0 \end{array} &
\begin{array}{c}Y_1\\ 1 \end{array} &
\begin{array}{c}Y_2\\ 2 \end{array} &
\begin{array}{c}Y_3\\ 3 \end{array} &
\begin{array}{c}Y_4\\ 4 \end{array} &
\begin{array}{c}Y_5\\ 5 \end{array}\\ \cline{2-8}
& \hfill 0 &
\tablebox{0}{0}{0}{} &
\tablebox{0}{0}{0}{} &
\tablebox{0}{0}{0}{} &
\tablebox{0}{0}{0}{} &
\tablebox{0}{0}{0}{} &
\tablebox{0}{0}{0}{} \\ \cline{2-8}
& X_1\hfill 1 &
\tablebox{0}{0}{0}{} &
\tablebox{0}{0}{0}{} &
\tablebox{0}{0}{0}{} &
\tablebox{0}{0}{0}{} &
\tablebox{0}{0}{0}{} &
\tablebox{0}{0}{0}{} \\ \cline{2-8}
& X_2\hfill 2 &
\tablebox{0}{0}{0}{} &
\tablebox{0}{0}{0}{} &
\tablebox{0}{0}{0}{} &
\tablebox{0}{0}{0}{} &
\tablebox{0}{0}{0}{} &
\tablebox{0}{0}{0}{} \\ \cline{2-8}
& X_3\hfill 3 &
\tablebox{0}{0}{0}{} &
\tablebox{0}{0}{0}{} &
\tablebox{0}{0}{0}{} &
\tablebox{0}{0}{0}{} &
\tablebox{0}{0}{0}{} &
\tablebox{0}{0}{0}{} \\ \cline{2-8}
i \hspace*{8pt} & X_4\hfill 4 &
\tablebox{0}{0}{0}{} &
\tablebox{0}{0}{0}{} &
\tablebox{0}{0}{0}{} &
\tablebox{0}{0}{0}{} &
\tablebox{0}{0}{0}{} &
\tablebox{0}{0}{0}{} \\ \cline{2-8}
& X_5\hfill 5 &
\tablebox{0}{0}{0}{} &
\tablebox{0}{0}{0}{} &
\tablebox{0}{0}{0}{} &
\tablebox{0}{0}{0}{} &
\tablebox{0}{0}{0}{} &
\tablebox{0}{0}{0}{} \\ \cline{2-8}
& X_6\hfill 6 &
\tablebox{0}{0}{0}{} &
\tablebox{0}{0}{0}{} &
\tablebox{0}{0}{0}{} &
\tablebox{0}{0}{0}{} &
\tablebox{0}{0}{0}{} &
\tablebox{0}{0}{0}{} \\ \cline{2-8}
& X_7\hfill 7 &
\tablebox{0}{0}{0}{} &
\tablebox{0}{0}{0}{} &
\tablebox{0}{0}{0}{} &
\tablebox{0}{0}{0}{} &
\tablebox{0}{0}{0}{} &
\tablebox{0}{0}{0}{} \\ \cline{2-8}
\end{array}
\end{displaymath}
\caption{Table template for dynamic programming solution
of BLCS instance}
\label{figure:emptyBLCS}
\end{figure}

\clearpage




\problem{30}
Assume that you are working at an animal shelter.
%
Suppose that there are $n$ puppies in the shelter and each of them has a specific hunger level $h_i$. 
%
Today you are given  $m$ biscuits, each of them has a size $b_j$.
%
Your job  today is to feed and satisfy as many hungry puppies as you can according to the biscuits that you have.
%
If a puppy has hunger level $h_i$, it can only be satisfied by feeding a biscuit of size $b_j \ge h_i$. 
%
That is, the biscuit size should be greater than or equal to the hunger level of the puppy to satisfy it.
%
Of course you cannot give the same biscuit to two puppies and each puppy can only get one biscuit.

%

{\bfseries
	\begin{enumerate}
		
		\begin{SLOPPY}
			
			\item
			State the {\normalfont{\textsc{Puppy Feeding}}} problem
			in the formal instance/solution format
			that we use in class.
			
		\end{SLOPPY}
		
		\item
		Design an efficient algorithm
		to solve the {\normalfont{\textsc{Puppy Feeding}}} problem
		and give the corresponding pseudocode.
		Argue that your algorithm returns an optimal solution.
		
		\item
		Give the $\Theta$ asymptotic worst-case time complexity
		for your algorithm.
		
	\end{enumerate}
}


\solution

% PUT YOUR SOLUTION HERE


\problembreak


\clearpage

\problem{20}
Let $G=(V,E)$ be an undirected (not necessarily connected) graph
with each edge colored either blue or red.
We are interested in whether the vertices in $V$
can be assigned to two categories: $C_1$ and $C_2$.
We take a blue edge $(u,v)$ to mean
that $u$ and $v$ should be in the same category,
while we take a red edge $(u,v)$ to mean
that $u$ and $v$ should be in different categories.
$G$ is \textit{assignable}
if there is an assignment of vertices to $C_1$ and $C_2$
such that the meaning of the edges is correct;
$G$ is \textit{unassignable} otherwise.

{\bfseries
\begin{enumerate}

\item
Draw an example of a graph $G$ with six vertices, four blue edges,
and four red edges
that is assignable.
Give an assignment of the vertices to $C_1$ and $C_2$
that demonstrates that $G$ is assignable.

\item
Draw an example of a graph $G$ with six vertices, four blue edges,
and four red edges
that is unassignable.
Explain why $G$ is not assignable.

\item
Formally state the problem {\normalfont{\textsc{Assignability}}}
that answers whether an edge-colored graph is assignable.

\item
Using BFS or DFS,
design an algorithm {\normalfont{\textsc{Assignable}}}
to solve {\normalfont{\textsc{Assignability}}}.
Write {\normalfont{\textsc{Assignable}}} in pseudocode.
Explain the major changes
required to either BFS or DFS
to achieve your {\normalfont{\textsc{Assignable}}} algorithm.

\item
Analyze the worst-case time complexity of {\normalfont{\textsc{Assignable}}}.

\end{enumerate}
}

\solution

% PUT YOUR SOLUTION HERE

\problembreak

\end{document}
